% !TeX root = ../papers/nc-p3.tex
\section{階数14の全成分上の四次中心}
\label{sec:rank-fourteen-center}
\begin{theorem}\label{thm:rank-fourteen-center}
標数零で、二本の階数14成分の全ての通常点で四次中心の文字次元は
$(3,2,0,0)$ である。有理な倍率座標 $t$ の上で四つの既約閉点
（次数 $8,8,20,34$）を除けば、関数体の消去で得た五つの中心元が各ファイバーの基底を与える。
この四閉点の上の全280幾何点は標準Koszul左右AS正則である。
\end{theorem}
\begin{proof}
$K=\mathbb Q(t)(i,Y)$、$i^2=-1,\ Y^2=C=2(t-1)(t+1)(t+3)$ と置く。
前節の大域表示により、二位置移動の二根と $i$ の二値を同じ有理係数の計算で扱える。
\path{computations/rank_fourteen_function_field.py} は $K$ の元を
$(A+Bi+DY+EiY)/q$、$A,B,D,E,q\in\mathbb Z[t]$ と表示し、共通因子を割る。
分子の $\mathbb Q(t)$ へのノルムは
\[
 \{A^2-B^2-C(D^2-E^2)\}^2+4(AB-CDE)^2 .
\]
この計算では数値補間を用いない。

始点・長さ・文字ごとに、直前の商基底の左に一文字を付け、
始端の二次関係に全ての残る語を付けて行消去する。この規則は全ての関係の挿入を含む。
自由列の語への射影を長さ5まで記録する。
\path{computations/verify_rank_fourteen_center.py} はこの既約分母を全て因数分解する。
通常倍率の定義で逆にした因子を除くと、残る極の第一座標の式はただ一つである。
\[
 t^8+8t^7+60t^6+248t^5+518t^4+568t^3+828t^2+1096t+769 .
\]
その外では各射影は定義され、関係を零にし、選んだ語を単位基底に送る。
四次・五次の道空間の次元は全点で既知なので、この全射はファイバーごとに同型となる。
三次以下でも同じ結論が成り立つ（固定曲線で標準の三次交叉を持つ）。
従って掃き出しの途中のピボットが消えても、最終射影が正則な点は取り除く必要がない。

各文字 $h$ の中心条件を、四始点の四次の語 $v_p$ に対する
$v_p x_j-x_jv_{p+1}=0$ の五次座標で記す。各行列は224行である。
独立な行を順に選び、以前のピボット列を零にした後に新たな列を一つ選ぶ。
選んだ行・列の順序では、ピボットの積はこの小行列式そのものである。
その積を体内で簡約してからノルムの分子を因数分解すると次の表を得る。
最終列には各小行列式について通常倍率の単元を除いた相異なる既約因子の次数を記す。
\[
\begin{array}{c|r|r|l}
 h&\text{列数}&\text{階数}&\text{因子の次数}\\ \hline
 0&36&33&8,8,20\\
 1&40&38&8,34\\
 2&32&32&8\\
 3&32&32&8
\end{array}
\]
表の和集合は上記の極を含む四因子である。各行列の一般階数は特殊化での上界となり、
選んだ小行列式が可逆な各点では等号である。余った列を自由変数に取ると、Cramerの公式で
この開集合上に中心の基底が得られる。

残る四因子については、各既約式の単純な素体根を一つずつ取る。
定数項193と769の八次式について $(p,t)=(157,141),(101,89)$、
二十次式と三十四次式については $(101,45),(101,85)$ を用いる。
全て通常な還元で、$i^2=-1$ と $s^2=w(t)$ はともに相異なる二根を持つ。
各四組で同じ中心次元を得る。また四つの中心四次式で割った商の次元は、
全始点で命題\ref{prop:fixed-residual-regularity}の次数0から13の表に一致する。
単根を持ち上げ、四組の拡大を全て比較したので、既約な $t$ の全共役が検査される。
同命題の定階数と中心持上げによって、これらの全点で全次数の正則性を得る。
\end{proof}
